% GDZ PPN236018647,native501;HKUST502;original diary5v.
% Complete entries corresponding to printed72-75 and77-80.
\clearpage\SourceLeaf{501}\NoSourceHead
\begingroup\fontsize{10}{15}\selectfont\itshape
\newcommand{\DiaryStrike}[1]{\tikz[baseline=(t.base)]{\node[inner sep=0pt,outer sep=0pt] (t) {\strut #1};\draw[line width=.6pt] ([yshift=.1ex]t.west)--([yshift=.1ex]t.east);}}
\noindent\tikz[baseline=-.4ex,x=1em,y=1em]{
\draw[line width=.3pt] (-.1,-1)--(.45,1.1);
\draw[line width=.3pt] (.12,-1)--(.64,1.12);
\draw[line width=.3pt] (.35,-.95)--(.84,1.05);
\draw[line width=.7pt] (.02,.85)--(.99,.55);
\draw[line width=.7pt] (-.01,.52)--(.96,.24);
\draw[line width=.7pt] (-.06,.2)--(.93,-.08);
}\quad Plani possibilitatem demonstravi.\hfill Jul. 28
\par\noindent\hfill Gotting.

\medskip
\noindent\underline{Quod Jul. 27. inscrips. errorem involvit: sed eo felicius}\\
\underline{nunc rem exhausimus, quoniam probare possumus nullum}\\
\underline{periodum esse posse numerum \DiaryStrike{fract} rationalem.}\hfill Aug 1

\medskip
\noindent\underline{Quomodo. periodorum numerum duplicando}\\
\hspace*{2em}\underline{signa adornare oporteat}

\medskip
\noindent\underline{Functionum primarum multitudinem per analysin}\\
\hspace*{2em}\underline{simplicissimam erui.}\hfill Aug 26

\medskip
\noindent\underline{Demonstratum, viaque ad multa maiora per intro\={d}.}\\
\hspace*{2em}\underline{Modulorum multiplicium strata}\hfill Aug. 31.

\medskip
\noindent Aug. 1. \underline{generalius ad quosvis modulos adaptatur}
\par\noindent\hfill Sept. 4.

\medskip
\noindent\tikz[baseline=-.3ex,x=1em,y=1em]{
\draw[line width=.3pt] (-.15,-.9)--(.55,1);
\draw[line width=.3pt] (.13,-1)--(.93,1.12);
\draw[line width=.3pt] (.03,.43)--(.5,-.85);
\draw[line width=.3pt] (.35,.6)--(.8,-.89);
\draw[line width=.7pt] (-.02,.35)--(.91,.3);
\draw[line width=.7pt] (-.12,-.35)--(.96,-.34);
}\quad\underline{Principia detexi, ad quae congruentiarum secundum}\\
\hspace*{2em}\underline{modulos multiplices resolutio ad congruentias secund}\\
\hspace*{2em}\underline{modulum linearem reducitur}\hfill Sept. 9

\medskip
\noindent\tikz[baseline=-.3ex,x=1em,y=1em]{
\draw[line width=.25pt] (-.05,-.85)--(.25,.95);
\draw[line width=.4pt] (.24,-1)--(.55,1.1);
\draw[line width=.25pt] (.48,-1.4)--(.72,.85);
\draw[line width=.5pt] (-.02,.6)--(1,.48);
\draw[line width=.8pt] (.22,.12)--(.87,.08);
\draw[line width=.8pt] (.06,-.22)--(1.02,-.37);
}\quad\underline{Aequationes habere radices imaginarias}\\
\hspace*{2em}\underline{methodo genuina demonstratum.}\hfill Bruns. Oct.
\par\noindent\hspace*{3em}Prom. in dissert. pecul. Mense Aug. 1799:
\endgroup
